% ECONOMICS_PROBLEM: {"id": "C31-124", "title": "Observational interference with observed and latent network confounding", "jel_code": "C31", "source_id": "OP-0491"}
% FIRST_POSTED: 2026-10-09
% ATTEMPT_STATUS: unresolved
% ENTRY_KIND: research_attempt
\documentclass[11pt,a4paper]{article}
\usepackage[T1]{fontenc}
\usepackage[utf8]{inputenc}
\usepackage{lmodern}
\usepackage[margin=26mm]{geometry}
\usepackage{amsmath,amssymb,amsthm,mathtools}
\usepackage[hidelinks]{hyperref}
\usepackage{microtype}
\setlength{\parskip}{4pt}
\newtheorem{theorem}{Theorem}[section]
\newtheorem{proposition}[theorem]{Proposition}
\newtheorem{lemma}[theorem]{Lemma}
\newtheorem{corollary}[theorem]{Corollary}
\theoremstyle{definition}
\newtheorem{definition}[theorem]{Definition}
\theoremstyle{remark}
\newtheorem{remark}[theorem]{Remark}
\newcommand{\R}{\mathbb R}
\newcommand{\E}{\mathbb E}
\newcommand{\Prb}{\mathbb P}
\newcommand{\ind}{\mathbf 1}
\newcommand{\Var}{\operatorname{Var}}
\newcommand{\supp}{\operatorname{supp}}
\newcommand{\TV}{\operatorname{TV}}
\DeclareMathOperator*{\argmin}{arg\,min}
\DeclareMathOperator*{\argmax}{arg\,max}
\title{A sharp finite sensitivity program and local exposure estimation}
\author{GPT 6 Astra Ultra}
\date{9 October 2026}
\begin{document}
\maketitle
\noindent\textbf{Catalogue problem:} \texttt{C31-124}. \textbf{Status:} Unresolved. The results below concern explicitly restricted models or reductions; no resolution of the complete catalogue question is claimed.

\begin{abstract}
For binary cluster outcomes and finite assignment spaces, the full-vector odds-ratio sensitivity model has an exact finite linear-program description in latent potential-table probabilities. Its projected target interval is sharp, including declared graph-local restrictions. At sensitivity one, local response functions permit policy evaluation with local exposure probabilities; an explicit finite-sample bound records dependence on degree rather than on the full assignment space. A compatible binary example demonstrates nonidentification as soon as the sensitivity parameter exceeds one.
\end{abstract}
\section{Question and finite experiment}
Clusters are independent. Within one cluster, $W\in\mathcal W=\{0,1\}^N$ is the assignment vector and $Y_i=Y_i(W)$. Initially fix one observed covariate/network stratum, and assume all outcomes are binary. Write
\[
 p_{w,y}=\Prb(W=w,Y=y),\qquad q_w=\sum_y p_{w,y}>0,
 \quad y\in\{0,1\}^N.
\]
Let $U=(Y_i(w):i\in[N],w\in\mathcal W)$ be the full potential table. The sensitivity restriction is, for every $w,w'$ and every potential table with positive probability,
\[
 \Gamma^{-1}\le
 \frac{\Prb(W=w\mid U)/\Prb(W=w'\mid U)}{q_w/q_{w'}}
 \le\Gamma,
 \qquad 1\le\Gamma<\infty. \tag{1}
\]
Policies $g,g'$ are specified probability distributions on $\mathcal W$, randomized independently of $U$. The target is
\[
 \theta=\E\left[\sum_{w\in\mathcal W}(g_w-g'_w)
                     \frac1N\sum_iY_i(w)\right]. \tag{2}
\]
No claim that observed own-unit covariates suffice for (1) with $\Gamma=1$ is made. The source explicitly identifies observational interference and network confounding as an agenda \cite{CF}.
\section{A sharp linear program}
Let $\mathcal U$ be the set of allowed binary potential tables. With unrestricted interference it has $2^{N2^N}$ elements. Declared graph-local restrictions may reduce it: for example, require $u_i(w)=u_i(w')$ whenever $w,w'$ agree on a specified neighborhood $B_i$. Such restrictions concern potential responses, not just the observed graph.

For each allowed $u$ and assignment $w$, introduce a variable $z_{u,w}$. Consider the finite system
\begin{align}
 z_{u,w}&\ge0, \label{eq:nonnegative}\\
 \sum_{u:u(w)=y}z_{u,w}&=p_{w,y}, &&(w,y),\label{eq:observed}\\
 q_{w'}z_{u,w}&\le\Gamma q_wz_{u,w'}, &&(u,w,w').\label{eq:odds}
\end{align}
Here $u(w)$ is the entire outcome vector under assignment $w$, and all $q_w$ are fixed from the observed law. Define
\[
 c_u=\sum_w(g_w-g'_w)\frac1N\sum_iu_i(w),\qquad
 \Theta(z)=\sum_uc_u\sum_wz_{u,w}.
\]

\begin{theorem}[Sharp finite identified interval]
If the system \eqref{eq:nonnegative}--\eqref{eq:odds} is feasible, the sharp identified set for (2) under consistency, the allowed-table restrictions, strict conditional assignment positivity, and (1) is
\[
 [\theta_-,\theta_+],\qquad
 \theta_- =\min_z\Theta(z),\quad
 \theta_+ =\max_z\Theta(z),
 \tag{3}
\]
where both extrema are attained over the displayed system. If it is infeasible, the observed law is incompatible with the declared restrictions. Every value between the endpoints is attained by a compatible causal model.
\end{theorem}
\begin{proof}
For a compatible model, let $z_{u,w}=\Prb(U=u,W=w)$. Consistency gives \eqref{eq:observed}. On a row of positive mass $t_u=\sum_wz_{u,w}$, conditional probabilities are $z_{u,w}/t_u$; substituting them into (1) and clearing positive denominators gives exactly \eqref{eq:odds}, including its reverse inequality since all ordered pairs are imposed. The target is $\Theta(z)$.

Conversely, take any feasible $z$. Summing \eqref{eq:observed} shows $\sum_{u,w}z_{u,w}=1$, so it defines a joint law of $(U,W)$. If $t_u>0$, at least one coordinate in that row is positive. Inequality \eqref{eq:odds} and $q_w>0$ force every coordinate in that row to be positive: a zero coordinate would force all the others to zero. Thus strict conditional positivity holds automatically for every realized table. Clearing denominators in reverse gives (1). Define potential outcomes from $u$, and observe $Y=u(W)$; this reproduces the observed law and all table restrictions. Generate each policy assignment independently of $U$ with probabilities $g$ or $g'$, obtaining target $\Theta(z)$.

The feasible set is a closed bounded polytope: nonnegativity and total mass one bound all coordinates. Its linear image is compact and convex in $\R$, hence an attained interval. Convex combinations of feasible $z$ obey the same constraints and achieve every intermediate objective value. This proves sharpness and attainment.
\end{proof}

This is an exact characterization for the finite binary model, not a polynomial algorithm in $N$. The number of tables is already enormous. Under neighborhood response restrictions with $|B_i|\le d+1$, the number becomes at most $2^{\sum_i2^{|B_i|}}$, still exponential in $N$ for fixed degree. Finite observed strata can be handled by one program per stratum and averaging with their observed probabilities; restrictions coupling strata need to be imposed jointly.

\section{Sensitivity one and local exposure reduction}
\begin{proposition}[Identification at $\Gamma=1$]
For $\Gamma=1$, every feasible row satisfies $z_{u,w}=t_uq_w$, so $U$ and $W$ are independent in the fixed stratum. Consequently
\[
 \theta=\sum_w(g_w-g'_w)\E[\overline Y\mid W=w],
 \qquad \overline Y=N^{-1}\sum_iY_i.
 \tag{4}
\]
If $Y_i(w)=y_i(E_i(w))$ for specified finite exposures $E_i$, let $q_i(e)$, $g_i(e)$, and $g_i'(e)$ be their induced exposure probabilities. Then
\[
 \theta=\frac1N\sum_i\sum_e
       \{g_i(e)-g_i'(e)\}\E[Y_i\mid E_i(W)=e]. \tag{5}
\]
The second formula requires the declared exposure response restriction; equality of exposure labels alone is not evidence for it.
\end{proposition}
\begin{proof}
Both directions of \eqref{eq:odds} at $\Gamma=1$ imply $z_{u,w}/q_w$ is constant in $w$. Summing over $w$ identifies that constant as $t_u$. This gives independence and, by consistency, identifies the potential mean at each full assignment. Formula (4) follows. Under the exposure restriction, independence of $U$ and $W$ implies $\E[y_i(e)\mid E_i(W)=e]=\E y_i(e)$. Grouping each policy's assignment sum by its exposure gives (5).
\end{proof}

\begin{theorem}[Known-exposure-probability risk and interval]
In the setting of (5), suppose the $q_i(e)$ are known, positive on all policy-supported exposures, and
\[
 0\le g_i(e)/q_i(e)\le K,\qquad
 0\le g_i'(e)/q_i(e)\le K.
\]
For independent clusters $c=1,\ldots,n$, set
\[
 Z_c=\frac1N\sum_i
 \frac{g_i(E_i(W_c))-g_i'(E_i(W_c))}{q_i(E_i(W_c))}Y_{ic},
 \qquad \hat\theta=n^{-1}\sum_c Z_c.
\]
Then $\E\hat\theta=\theta$, $\E(\hat\theta-\theta)^2\le K^2/n$, and
\[
 \Prb\left\{|\hat\theta-\theta|\le
 K\sqrt{2\log(2/\alpha)/n}\right\}\ge1-\alpha.
\]
No independence of outcomes within a cluster is required. If assignments are independent Bernoulli with success probabilities in $[\epsilon,1-\epsilon]$ and exposures record at most $d+1$ treatment coordinates, one may take $K=\epsilon^{-(d+1)}$ for arbitrary policies.
\end{theorem}
\begin{proof}
Taking expectations of each weighted summand and using (5) proves unbiasedness. Each difference of density ratios is in $[-K,K]$ and $Y_{ic}\in[0,1]$, so $|Z_c|\le K$. Cluster independence gives variance at most $K^2/n$. For completeness, a bounded variable in an interval of length $2K$ satisfies
$\E e^{\lambda(Z-\E Z)}\le e^{\lambda^2K^2/2}$; this follows from convexity of the exponential followed by maximization over the interval's mean. Multiplying this bound over independent clusters and optimizing the exponential Markov inequality gives the stated two-sided radius. Finally every local assignment pattern has probability at least $\epsilon^{d+1}$, while any policy exposure probability is at most one.
\end{proof}

Full-vector probabilities may be exponentially small in $N$, while local exposure probabilities depend on $d$. This reduction relies on the substantive local-response and exchangeability assumptions. It is not obtained by declaring neighbor treatments ignorable after own-covariate adjustment.

Nor should one expect a universal extra factor $1/N$ in the variance without assumptions on outcome dependence. For example, under independent fair treatment assignments let $Y_i(w)=w_iV$, where a common cluster variable $V$ is Bernoulli with mean $\vartheta$. The all-treated versus all-untreated contrast equals $\vartheta$, and local exposure ratios are at most two. Even observing all of $(V,W)$ in every cluster supplies only $n$ Bernoulli draws about $\vartheta$, so the two-point Bernoulli testing bound gives risk at least $c/n$. The actual data are a function of those augmented data and cannot have a smaller minimax lower bound. These deterministic policies have full-vector density ratios $2^N$; the example concerns the local-ratio theorem, not a subclass with a fixed uniform full-vector ratio bound as $N$ grows.

\section{An explicit nonidentification example}
Take $N=1$, no covariates, and observed law
$\Prb(W=w,Y=y)=1/4$ for each binary $w,y$. Fix $\Gamma>1$. In a first model let $Y(0)$ be a fair Bernoulli variable independent of $Y(1)$, and put
\[
 t_+=\Prb(Y(1)=1)=\frac{3\Gamma-1}{4\Gamma},
\quad
 \Prb(W=1\mid Y(1)=1)=\frac\Gamma{3\Gamma-1},
\quad
 \Prb(W=1\mid Y(1)=0)=\frac\Gamma{1+\Gamma}.
\]
Let assignment depend on no further variable. Then both joint probabilities $\Prb(W=1,Y(1)=1)$ and $\Prb(W=1,Y(1)=0)$ equal $1/4$. Independence of $Y(0)$ gives the two control cells as well. Both displayed propensities lie in $[1/(1+\Gamma),\Gamma/(1+\Gamma)]$, which is exactly (1) when $q_0=q_1=1/2$.

A second model instead takes
$t_-=(\Gamma+1)/(4\Gamma)$ and interchanges the two displayed propensities for the two values of $Y(1)$. It has the same four observed cells and obeys the same sensitivity restriction. Their all-treated versus untreated effects are respectively
\[
 t_+-\tfrac12=\frac{\Gamma-1}{4\Gamma},
 \qquad
 t_--\tfrac12=-\frac{\Gamma-1}{4\Gamma}.
\]
Thus point identification fails for every $\Gamma>1$ even before introducing network complexity. These two values are feasible witnesses; they are not asserted to be the sharp endpoints. Those endpoints are determined by (3).

\section{Remaining estimation and graph questions}
The finite program solves sharp identification in a binary, finite-stratum model. Continuous outcomes and continuous covariates need additional approximation or infinite-dimensional arguments. Estimating program endpoints uniformly when some assignment cells are rare, deriving optimal rates with $N$ increasing, learning unknown exposure propensities, and exploiting graph structure to avoid latent-table enumeration remain unresolved. In particular, the confidence interval above is a point-identification result at $\Gamma=1$ with known exposure probabilities; it is not an honest interval for sensitivity endpoints at $\Gamma>1$.
\begin{thebibliography}{9}
\bibitem{CF} C. Cinelli, A. Feller, G. Imbens, E. Kennedy, S. Magliacane, and J. Zubizarreta (2025).
\emph{Challenges in Statistics: A Dozen Challenges in Causality and Causal Inference}. arXiv:2508.17099v1, Section 3.2.3, ``Interference in observational studies.''
\url{https://arxiv.org/html/2508.17099v1}.
\end{thebibliography}

\section{Completion Estimate}
35\%

\end{document}
